\section{Discussion and repair program}
\label{sec:discussion}

The structural decomposition remains useful despite the failed completeness
claim.  It identifies exactly where samples interact and provides a compact
description of infinitesimal compatibility.  The counterexample identifies the
minimum semantic information that any exact extension must add.

\subsection{Potential message enrichments}

The following are research directions rather than established polynomial-time
repairs.

\paragraph{Semialgebraic projection messages.}
Using the exact set in \cref{eq:zeroth-order-message} gives perfect semantics.
The unresolved issue is whether this special family has polynomial descriptive
complexity and supports polynomial-time elimination.  General quantifier
elimination does not supply the desired bound.

\paragraph{Affine or homogenized first-order messages.}
Augmenting a derivative matrix with constant terms prevents the immediate loss
of \(y\) in the audit family and is exact for affine constraint systems.  For
nonlinear systems, an affine row space still does not generally determine the
zero set or its projection.

\paragraph{Higher jets.}
Higher-order derivative data can retain more of a bounded-degree polynomial.
A successful method would need a uniform order, a basis-independent merge
operation, correct treatment of inequalities and real components, and a
polynomial representation bound.  None follows from the first-order separator
theorem.

\paragraph{Ideals, real radicals, and elimination ideals.}
An algebraic message could retain an elimination ideal rather than a cotangent
subspace.  Ordinary elimination ideals describe Zariski closures of complex
projections; exact real feasibility can additionally require inequalities and
real-radical information.  Gr\"obner or quantifier-elimination representations
may grow substantially.  The special block structure would have to yield a new
bound.

\subsection{Priority order for the target theorem}

Further work should proceed in the following dependency order:
\begin{enumerate}
  \item propose a zeroth-order message representation that distinguishes the
  audit family;
  \item prove its exact semantic invariant for leaf, merge, contraction, and
  root operations;
  \item attack it with small disconnected, singular, and incompatible-target
  instances;
  \item prove a polynomial representation-size theorem;
  \item derive a Turing bit-complexity bound;
  \item formalize the model-matching ETR reduction; and
  \item only then instantiate the conditional collapse theorem.
\end{enumerate}

This ordering prevents an arithmetic speedup for an auxiliary linear problem
from being mistaken for a complexity classification of the nonlinear decision
language.

\subsection{A publishable unconditional core}

Independently of the open collapse target, the following package is
unconditional: exact ReLU polynomialization including the boundary; the
shared-base tensor decomposition; the relative K\"ahler decomposition; the
block-angular Jacobian; the intrinsic linear CMP factorization; and the
in-scope impossibility theorem for cotangent-only exact decoding.  A future
positive algorithm can build on this audited base.
